% Part: proof-theory
% Chapter: proof-search
% Section: tableaux

\documentclass[../../../include/open-logic-section]{subfiles}

\begin{document}

\olfileid[hi]{pt}{ps}{tab}

\olsection{सार्थक वृक्ष}

सार्थक-वृक्ष प्रणालियाँ अनुक्रम कलनों से निकटतः संबंधित ऐसी !!{proof}
प्रणालियाँ हैं जो हाथ से और सॉफ्टवेयर में कार्यान्वित—दोनों प्रकार की
प्रमाण-खोज के लिए विशेष रूप से उपयुक्त हैं। अनुक्रमों के वृक्ष बनाने के बजाय
सार्थक-वृक्षों में !!a{proof}, !!{formula}s का वृक्ष होता है। किंतु प्राकृतिक
निगमन के विपरीत इनके नियम अनुक्रम कलन के नियमों का प्रतिबिंब हैं। सार्थक-वृक्ष
!!{proof} मूलतः किसी निदर्श की प्रत्यक्ष खोज है। इसलिए इन्हें कभी-कभी
\emph{अर्थविज्ञानपरक सार्थक-वृक्ष} या \emph{सत्य-वृक्ष} भी कहते हैं।

\emph{चिह्नित} सार्थक वृक्ष चिह्नित !!{formula}s का वृक्ष है, जिनका रूप
$\sFmla{\True}{!A}$ अथवा $\sFmla{\False}{!A}$ हो सकता है। वृक्ष नीचे की ओर
बढ़ते हैं। वे उन !!{formula}s से आरंभ होते हैं जो बताते हैं कि हम क्या
!!{prove} करना चाहते हैं। उदाहरणतः अनुक्रम $!A, !B \Sequent !C, !D$ को
!!{prove} करने के लिए वृक्ष $\sFmla{\True}{!A}, \sFmla{\True}{!B},
\sFmla{\False}{!C}, \sFmla{\True}{!D}$ से आरंभ होगा। ऐसी !!a{structure} जो
$!A$ और $!B$ को सत्य तथा $!C$ और $!D$ को असत्य बनाती है, यह दिखाने वाली
!!a{structure} है कि $!A, !B \Sequent !C, !D$ वैध नहीं है।

सार्थक वृक्ष के किसी पर्ण (सबसे निचले) आसंधि को शाखा-विस्तार नियमों द्वारा
बढ़ाया जा सकता है। ये नियम उसी शाखा में अतिरिक्त आसंधियाँ जोड़ते हैं या नीचे
दो सहोदर आसंधियाँ जोड़कर शाखा को दो नई शाखाओं में बाँटते हैं। नियमों को उसी
शाखा में उसके ऊपर स्थित किसी भी चिह्नित !!{formula} पर लगाया जा सकता है।
चिरसम्मत सार्थक-वृक्ष प्रणाली \Log{Tc} के शाखा-विस्तार नियम
\olref{tab:Tc} में दिए गए हैं।

\subfile{rules-Tc}

जैसा कि आप देख सकते हैं, ये नियम केवल \Log{G3c} के उलटे किए हुए नियम हैं;
$\True$ और $\False$ चिह्न क्रमशः बताते हैं कि !!a{formula} अनुक्रम के बाएँ या
दाएँ पक्ष पर प्रमुख है।

सार्थक वृक्ष की कोई शाखा \emph{बंद} है यदि उसमें $\sFmla{\True}{!A}$ और
$\sFmla{\False}{!A}$ दोनों हों। यदि प्रत्येक शाखा बंद है, तो हम पूरे सार्थक
वृक्ष को बंद कहते हैं। उदाहरणतः $!A \lor (!B \land !C) \Sequent (!A \lor
!B) \land (!A \lor !C)$ के लिए, अर्थात् $\sFmla{\True}{!A \lor (!B \land
!C)}$ और $\sFmla{\False}{(!A \lor !B) \land (!A \lor !C)}$ से आरंभ होने
वाला, एक संवृतक सार्थक वृक्ष यह है:
\begin{oltableau}
  [\sFmla{\True}{\formula{A} \lor (\formula{B} \land \formula{C})},
    just=\TAss, checked
    [\sFmla{\False}{(\formula{A} \lor \formula{B}) \land
        (\formula{A} \lor \formula{C})}, just=\TAss, checked
      [\sFmla{\True}{\formula{A}}, just={\TRule{\True}{\lor}[1]}
        [\sFmla{\False}{\formula{A} \lor \formula{B}},
          just={\TRule{\False}{\land}[2]}, checked
          [\sFmla{\False}{\formula{A}},
            just={\TRule{\False}{\lor}[4]}
            [\sFmla{\False}{\formula{B}},
              just={\TRule{\False}{\lor}[4]}, close]
          ]
        ]
        [\sFmla{\False}{\formula{A} \lor \formula{C}},
          just={\TRule{\False}{\land}[2]}, checked
          [\sFmla{\False}{\formula{A}},
            just={\TRule{\False}{\lor}[4]}
            [\sFmla{\False}{\formula{C}},
              just={\TRule{\False}{\lor}[4]}, close]
          ]
        ]
      ]
      [\sFmla{\True}{\formula{B} \land \formula{C}},
        just={\TRule{\True}{\lor}[1]}, checked
        [\sFmla{\False}{\formula{A} \lor \formula{B}},
          just={\TRule{\False}{\land}[2]}, checked
          [\sFmla{\True}{\formula{B}},
            just={\TRule{\True}{\land}[3]}, move by=2
            [\sFmla{\True}{\formula{C}},
              just={\TRule{\True}{\land}[3]}
              [\sFmla{\False}{\formula{A}},
                just={\TRule{\False}{\lor}[4]}
                [\sFmla{\False}{\formula{B}},
                  just={\TRule{\False}{\lor}[4]},close
                ]
              ]
            ]
          ]
        ]
        [\sFmla{\False}{\formula{A} \lor \formula{C}},
          just={\TRule{\False}{\land}[2]}, checked
          [\sFmla{\True}{\formula{B}},
            just={\TRule{\True}{\land}[3]}, move by=2
              [\sFmla{\True}{\formula{C}},
                just={\TRule{\True}{\land}[3]}
                [\sFmla{\False}{\formula{A}},
                  just={\TRule{\False}{\lor}[4]}
                  [\sFmla{\False}{\formula{C}},
                  just={\TRule{\False}{\lor}[4]},close
                ]
              ]
            ]
          ]
        ]
      ]
    ]
  ]
\end{oltableau}

सही के चिह्न बताते हैं कि संबंधित नियम को !!a{formula} के नीचे की सभी शाखाओं
में उस पर लगाया जा चुका है। $\otimes$ चिह्न बताता है कि शाखा बंद है। नियमों
के बाद की पंक्ति-संख्याएँ बताती हैं कि नियम किन चिह्नित !!{formula}s पर लगाया
गया है (अर्थात् उसी शाखा की बताई गई पंक्ति के !!{formula} पर)।

सार्थक वृक्षों में प्रमाण-खोज सरल है। हम सार्थक वृक्ष को चरणों में उत्पन्न
करते हैं। चरण~$0$ में आरंभिक !!{formula}s को ऊपर लिख दें। चरण~$n+1$ में
प्रत्येक पर्ण आसंधि के लिए यह करें: अभी तक सही का चिह्न न लगे हुए पहले (सबसे
ऊपरी) चिह्नित सूत्र को खोजें। उस सूत्र पर शाखा-विस्तार नियम लगाएँ। जिस प्रत्येक
शाखा में वह चिह्नित सूत्र है, उसमें नियम लग जाने के बाद उस पर सही का चिह्न लगा
दें। (ऊपर का उदाहरण इसी कलन-विधि से उत्पन्न हुआ है।)

$\sFmla{\True}{\lforall}$ और $\sFmla{\False}{\lexists}$ रूप के चिह्नित
!!{formula}s पर कभी सही का चिह्न नहीं लगाया जाता, इसलिए उनका विशेष उपचार
करना पड़ता है। यदि वे उपस्थित हों, तो हमें सुनिश्चित करना चाहिए कि अंततः सभी
लागू पदों के लिए $\TRule{\True}{\forall}$ और
$\TRule{\False}{\lexists}$ नियम लगाए जाएँ। उदाहरणतः प्रत्येक चरण में प्रत्येक
शाखा के ऐसे सभी सूत्रों पर इन्हें लगाकर इसकी गारंटी दी जा सकती है (इस रूप का
न होने वाले सर्वोच्च चिह्नित !!{formula} से निपटने के बाद)। प्रयुक्त पद~$t$
ऐसा पहला पद है जिसे शाखा में पहले से आने वाले !!{constant}s (या कोई
!!{constant} न आए तो $\Obj c_0$) और आरंभिक !!{formula}s में आने वाले
!!{function}s से बनाया जा सकता है, और जिसका संगत चिह्नित !!{formula}
$\sFmla{\True}{!B(t)}$ अथवा~$\sFmla{\False}{!B(t)}$ शाखा में पहले से नहीं है।

यदि यह प्रमाण-खोज बंद सार्थक वृक्ष उत्पन्न नहीं करती, तो या तो उसमें ऐसी
सीमित खुली शाखा है जहाँ सभी गैर-परमाण्विक !!{formula}s पर सही का चिह्न लग चुका
है, या खोज ऐसा अनंत, सीमित-शाखी वृक्ष उत्पन्न करती है जिसमें कोई अनंत शाखा
होनी ही चाहिए। \olref[cpl]{sec} की ही तरह हम ऐसी !!a{structure}~$\Struct{M}$
परिभाषित कर सकते हैं कि शाखा में $\sFmla{\True}{!A}$ आने पर $\Sat{M}{!A}$,
और $\sFmla{\False}{!A}$ आने पर $\Sat/{M}{!A}$ हो। ($\Theta =
\Setabs{!A}{\sFmla{\True}{!A} \text{ शाखा में}}$ और $\Xi =
\Setabs{!A}{\sFmla{\False}{!A} \text{ शाखा में}}$ लें।)

इस प्रमाण-खोज विधि से उत्पन्न सार्थक वृक्षों का प्रत्येक आसंधि को एक अनुक्रम
देकर आसानी से \Log{G3c} के प्रमाणों में अनुवाद किया जा सकता है। यदि आरंभिक
सूत्र अनुक्रम $\Gamma \Sequent \Delta$ के अनुरूप हैं, तो प्रत्येक आरंभिक सूत्र
को $\Gamma \Sequent \Delta$ से चिह्नित करें। यदि किसी शाखा को चिह्नित
!!{formula} $\sFmla{\True}{!A}$ पर शाखा-विस्तार नियम लगाकर बढ़ाया जाता है, तो
पर्ण को $!A, \Pi \Sequent \Lambda$ से चिह्नित करें; और यदि नियम चिह्नित
!!{formula} $\sFmla{\False}{!A}$ पर लगा है, तो पर्ण को $\Pi \Sequent
\Lambda, !A$ से चिह्नित करें। जहाँ सार्थक वृक्ष $\TRule{\True}{\star}$ नियम
लगाता है, वहाँ अनुक्रम पर \LeftR{\star} नियम ($!A$ को प्रमुख रखकर) लगाएँ;
जहाँ वह $\TRule{\False}{\star}$ नियम लगाता है, वहाँ \RightR{\star} नियम
लगाएँ। नियम के आधार-वाक्य सार्थक वृक्ष में जोड़े गए संगत आसंधियों को चिह्नित
करने वाले अनुक्रम हैं। जब कोई सार्थक वृक्ष $\sFmla{\True}{!A}$ और
$\sFmla{\False}{!A}$ से बंद होता है, तब अंतिम पर्ण आसंधि का चिह्न $!A, \Pi
\Sequent \Lambda, !A$ रूप का अनुक्रम होगा। सार्थक वृक्ष की कुछ क्रमिक
आसंधियों को एक ही अनुक्रम मिलेगा; उनकी प्रतियाँ हटा दें। उदाहरणतः ऊपर का
सार्थक वृक्ष इस !!{proof} में अनूदित होता है:
\[\small
\Axiom$!A \fCenter !A, !B$
\RightLabel{\RightR{\lor}}
\UnaryInf$!A \fCenter !A \lor !B$
\Axiom$!A \fCenter !A, !C$
\RightLabel{\RightR{\lor}}
\UnaryInf$!A \fCenter !A \lor !C$
\RightLabel{\RightR{\land}}
\BinaryInf$!A \fCenter (!A \lor !B) \land (!A \lor !C)$
\Axiom$!B, !C \fCenter !A, !B$
\RightLabel{\RightR{\lor}}
\UnaryInf$!B, !C \fCenter !A \lor !B$
\RightLabel{\LeftR{\land}}
\UnaryInf$!B \land !C \fCenter !A \lor !B$
\Axiom$!B, !C \fCenter !A, !C$
\RightLabel{\RightR{\lor}}
\UnaryInf$!B, !C \fCenter !A \lor !C$
\RightLabel{\LeftR{\land}}
\UnaryInf$!B \land !C \fCenter !A \lor !C$
\RightLabel{\RightR{\land}}
\BinaryInf$!B \land !C \fCenter (!A \lor !B) \land
(!A \lor !C)$
\RightLabel{\LeftR{\lor}}
\BinaryInf$!A \lor (!B \land !C) \fCenter (!A \lor !B) \land
(!A \lor !C)$
\DisplayProof
\]

\end{document}
